\chapter{Preface}

This book has a simple ambition: a single, coherent chemistry book
covering everything from the first year of school to the end of the
bachelor's degree, written in English for readers anywhere in the world.

The style is concise and rigorous: courses built from definitions, examples,
propositions, methods and diagrams, with careful reasoning whenever it is
accessible at the given level, followed by carefully chosen exercises.
Results stated without full derivation are explicitly marked as admitted.
Full solutions to all exercises are collected at the end of the book.
Chemical equations are always balanced, and every number quoted for a
real substance comes from a reference table.

The course comes in four books. The first covers all twelve school years
in a single volume and never teaches the same notion twice: each idea is
taught once, in the year where it first fits, and a later chapter that
needs it opens with a short reminder, \emph{You already know}, before
going further. The other three books are the three years of university.
In the first book each school year is a part, split into chapters; each
university book is one year, split into chapters. The earlier the year, the younger the reader it is written
for: more illustrations, more detailed worked steps, and gentler
exercises.

\section*{How to read this book}

Statements are numbered within each chapter and color-coded: definitions in
blue, theorems in red, propositions and their relatives in orange, and
methods --- short recipes for standard tasks --- in green. Exercises are
graded from ($\star$) for direct applications to ($\star\star\star$) for
harder problems. Framed boxes hold what is not part of the argument: a
reminder of what an earlier chapter taught, how a procedure is carried
out in a laboratory, a page of history, or the hazards of a product.

\section*{Safety}

Chemistry is learnt with real substances, and some of them are dangerous.
The experiments described in this book are carried out in a school or
university laboratory, by or with a teacher, with the protective
equipment the laboratory provides. None of them is meant to be done at
home.

\section*{Contributing}

The source of this book lives in a public git repository:
\begin{center}
\url{https://github.com/bvirrion/one-chemistry-book}
\end{center}
Contributions --- new chapters, corrections, better explanations,
additional exercises --- are welcome; see \texttt{CONTRIBUTING.md} at the
root of the repository.
